\documentclass[11pt]{article}
\newcommand{\ArticleShort}{Synthesis}
\usepackage[T1]{fontenc}
\usepackage[utf8]{inputenc}
\usepackage[english]{babel}
\usepackage{newtxtext}
\usepackage{amsmath,amsthm,mathtools}
\usepackage{newtxmath}
\usepackage[letterpaper,margin=1in,headheight=15pt]{geometry}
\usepackage{microtype,setspace,booktabs,tabularx,longtable,array,enumitem,titlesec,fancyhdr}
\usepackage{xcolor}
\definecolor{ink}{HTML}{193244}
\usepackage{csquotes}
\usepackage[backend=biber,style=apa,natbib=true,doi=true,url=true,eprint=false]{biblatex}
\addbibresource{shared/references.bib}
\usepackage[unicode,hidelinks,bookmarksnumbered=true]{hyperref}
\usepackage{bookmark,xurl}
\setstretch{1.07}
\setlength{\parindent}{1.3em}
\setlength{\parskip}{0.2em}
\setlength{\emergencystretch}{2em}
\setlength{\bibitemsep}{0.45\baselineskip}
\setlength{\bibhang}{1.5em}
\renewcommand*{\bibfont}{\small}
\setcounter{biburllcpenalty}{7000}
\setcounter{biburlucpenalty}{7000}
\setcounter{biburlnumpenalty}{7000}
\setlist{topsep=0.4em,itemsep=0.25em,parsep=0pt}
\titleformat{\section}{\large\bfseries\color{ink}}{\thesection}{0.65em}{}
\titleformat{\subsection}{\normalsize\bfseries}{\thesubsection}{0.65em}{}
\titlespacing*{\section}{0pt}{1.2\baselineskip}{0.45\baselineskip}
\titlespacing*{\subsection}{0pt}{0.8\baselineskip}{0.3\baselineskip}
\pagestyle{fancy}
\fancyhf{}
\fancyhead[L]{\small\itshape \AE{}ther-flow and General Relativity}
\fancyhead[R]{\small\ArticleShort}
\fancyfoot[L]{\scriptsize Revised manuscript \textperiodcentered{} 9 September 2026}
\fancyfoot[R]{\small\thepage}
\renewcommand{\headrulewidth}{0.3pt}
\renewcommand{\footrulewidth}{0pt}
\newcommand{\Aether}{\AE{}ther}
\newcommand{\Aflow}{\AE{}ther-flow}
\newcommand{\TheoryName}{\Aflow{} Interpretation of Relativity}
\newcommand{\TheoryProgram}{\Aether{} / \Aflow{} framework}
\newcommand{\Sub}{\mathcal A}
\newcommand{\PhiSrc}{\Phi_{\mathrm{src}}}
\newcommand{\Order}{\mathcal S}
\newcommand{\Ph}{\Phi}
\newcommand{\Proj}{D}
\newcommand{\TT}{\mathrm{TT}}
\newcommand{\dd}{\mathrm d}
\newcommand{\R}{\mathbb R}
\newtheorem{definition}{Definition}
\newtheorem{proposition}{Proposition}
\theoremstyle{remark}
\newtheorem{remark}{Remark}
\newcommand{\PaperTitle}[3]{%
  \hypersetup{pdftitle={#1},pdfsubject={#2},pdfauthor={}}%
  \thispagestyle{plain}%
  \begin{center}
  {\small\scshape \AE{}ther-flow and General Relativity\par}
  \vspace{0.7em}
  {\LARGE\bfseries\color{ink}#1\par}
  \vspace{0.5em}
  \if\relax\detokenize{#2}\relax\else
  {\normalsize #2\par}
  \vspace{0.35em}
  \fi
  \end{center}
  \begin{abstract}\noindent #3\end{abstract}
  \vspace{0.35em}
}
\newcommand{\References}{\printbibliography[heading=bibintoc,title={References}]}

% Copyright footer added 10 September 2026.
\fancyfoot[L]{\scriptsize Revised manuscript \textperiodcentered{} 15 September 2026 \textperiodcentered{} Copyright \textcopyright{} 2026 Alexander Samuel Ricciardi \textperiodcentered{} AngryOwl}
\fancyfoot[R]{\small\thepage}
\fancypagestyle{plain}{%
  \fancyhf{}%
  \fancyfoot[L]{\scriptsize Revised manuscript \textperiodcentered{} 15 September 2026 \textperiodcentered{} Copyright \textcopyright{} 2026 Alexander Samuel Ricciardi \textperiodcentered{} AngryOwl}%
  \fancyfoot[R]{\small\thepage}%
  \renewcommand{\headrulewidth}{0pt}%
  \renewcommand{\footrulewidth}{0pt}%
}

\begin{document}
\PaperTitle{\Aflow{} as an Interpretation of General Relativity}{Source assumptions, observer geometry, and the explanatory gap}{The \Aether{} / \Aflow{} proposal seeks a physical account of spatial and temporal organization through intrinsic ordered motion in a proposed four-dimensional substrate, while retaining general relativity as the recovery target. Its present mathematical assumptions are limited to a smooth four-dimensional arena and an unresolved source-flow placeholder. General relativity is adopted separately as the quantitative benchmark, with observer congruences supplying an effective interpretation of flow. This article relates that distinction to a dual research objective: test a source explanation of relativity and evaluate whether an artificial-intelligence research system can make independently checkable progress. It specifies assumption control, recovery targets, scoped negative outcomes, and a proposed evaluation protocol. Fixed-assumption audits are separated from constructions that add new premises. The work establishes the boundaries and evidence requirements of this research case study, not a source derivation, a universal impossibility theorem, or measured model performance.}

\section{The proposal and its central distinction}
General relativity (GR) gives a quantitative account of spacetime geometry and its relation to matter. The \TheoryProgram{} asks a different, additional question: could the spatial and temporal relations used in that account arise from a source organization not identified at the outset with physical spacetime? The language of \Aether{} and \Aflow{} is intended to express this source-level possibility.

A possibility is not yet a model. To explain spacetime from another structure, one must define the source objects, state their laws or constraints, and demonstrate how the operational content of relativity follows. The present source specification has not completed these steps. Its effective model instead takes GR as an explicit starting assumption.

The distinction prevents two opposite mistakes. The first would be to announce that the ontology has derived Einsteinian gravity because its effective equations have been chosen to be Einstein's equations. The second would be to conclude that the source proposal must be false simply because a derivation has not been supplied. Neither conclusion follows from the stated assumptions. The framework provides an adopted comparison model and an unfinished explanatory proposal.

The research goal is to test that explanatory proposal and, separately, the ability of an AI-assisted system to produce reliable results about it. The evidence requirements below allow a valid scoped negative result to count as successful research.

This article brings those parts together. The technical companion papers give the action variation, perturbative mode count, operational limits, and congruence calculations. Their standard GR content should be read as inherited physics, not as a claim that the proposal has independently discovered those results.

\section{What is assumed about the source}
The source arena is denoted by $\Sub$ and assumed to be a smooth four-dimensional manifold. It is not identified with the effective spacetime manifold $M$. Its dimension and smooth structure are primitive assumptions, not products of a source derivation. No particular metric, physical causal order, clock, or matter content is supplied on $\Sub$.

The symbol $\PhiSrc$ is a placeholder for a possible source-order or evolution structure. Its mathematical type remains unresolved. It has not been specified as a one-parameter flow, a vector field, a stochastic process, or a generator. The usual properties of a mathematical flow therefore cannot be inferred from its name. In particular, no parameter has been identified as physical time.

Intrinsic ordered motion is the working physical hypothesis, not merely an alternative name for observer motion in GR. The proposed substrate is meant to generate the metric rather than move through a previously specified space. Observers belong to the same underlying reality; their accessible spatial and temporal relations would be local manifestations of its organization. This is the intended explanation, not a mechanism already constructed from $\Sub$ and $\PhiSrc$.

Experienced space and S-time retain their roles in that hypothesis. A local spatial slice expresses the intended account of three-dimensional measurement, while S-time links temporal experience to ordered change. The present assumptions do not define an embedding of that slice in $\Sub$, a monotone clock functional, or an explanation of subjective experience. They also do not authorize treating an unspecified source parameter as physical time. Foundations separates these conceptual commitments from the mathematical structures that a reconstruction must supply.

This restraint does not remove the original source question. It makes the proposed explanation evaluable in principle. A future source model must say what carries the order, what counts as a state or process, which transformations are allowed, and how a physical observer can obtain measurements from it. Until those relations exist, referring to $\Sub$ as underlying reality is a hypothesis about physical interpretation, not an experimentally established fact.

\section{A definite effective gravitational model}
\subsection{Fields and equations}
The effective arena is a time-oriented four-dimensional Lorentzian spacetime $(M,g)$ with signature $(-+++)$ and Levi-Civita connection. In length-valued covariant coordinates, $x^0=ct$. With a physical matter action $S_m$, the adopted model is
\begin{equation}
S_{\rm eff}[g,\psi]=\frac{c^3}{16\pi G}
\int_M\dd^4x\sqrt{-g}(R-2\Lambda)+S_m[g,\psi],
\label{eq:SeffLambda}
\end{equation}
where $G>0$, $\Lambda$ is constant, and $\psi$ denotes specified ordinary matter. The $\Lambda=0$ specialization is
\begin{equation}
S_{\rm eff}\big|_{\Lambda=0}=\frac{c^3}{16\pi G}
\int_M\dd^4x\sqrt{-g}R+S_m[g,\psi].
\label{eq:Seff}
\end{equation}
The stress tensor and bulk metric equation are
\begin{equation}
T_{\mu\nu}=-\frac{2c}{\sqrt{-g}}\frac{\delta S_m}{\delta g^{\mu\nu}},
\qquad G_{\mu\nu}+\Lambda g_{\mu\nu}=\frac{8\pi G}{c^4}T_{\mu\nu}.
\label{eq:Einstein}
\end{equation}
These are the ordinary GR dynamics with the action units stated explicitly. Compactly supported variations suffice for the bulk equation; specific boundary problems need the corresponding boundary terms and data.

No independent $\PhiSrc$ field appears in \eqref{eq:SeffLambda}. Thus there is no additional low-energy gravitational equation, coupling parameter, or source-flow mode to compare with the metric sector. The usual local gravitational degree-of-freedom count is retained \parencite{ADM1962}. That is a property of the chosen action, not evidence that an undefined microscopic theory has already reduced to two modes.

\subsection{Operational meaning}
For standard Maxwell light in vacuum geometric optics,
\begin{equation}
g_{\mu\nu}k^\mu k^\nu=0,
\label{eq:null}
\end{equation}
with null geodesic propagation. For an ideal timelike clock,
\begin{equation}
\dd\tau^2=-\frac1{c^2}g_{\mu\nu}\dd x^\mu\dd x^\nu.
\label{eq:tau}
\end{equation}
These formulas connect the metric to measurement \parencite{Carroll1997}. They require the ordinary matter laws used to model light and clocks; the phrase minimal coupling is not a proof that an arbitrary matter action has the same propagation properties.

The benchmark agrees exactly with the corresponding GR model only when matter, constants, allowed solutions, initial and boundary data, and measurement procedures match. A hidden source constraint selecting different solutions would change that content. Here no such constraint is specified. The statement of agreement is consequently conditional model identity, not an additional empirical success attributable to the ontology.

The proposed physical explanation aims to preserve this agreement without an additional direct flow force on matter. The intended mechanism would recover the common metric and its coupling to measuring systems. Adopting metric-only coupling and exact GR dynamics already gives the reference model, but does not establish that a deeper source produces it. Nor does metric-only matter coupling alone exclude observable preferred structure in a candidate with different gravitational dynamics or restrictions. Those dependencies belong in the reconstruction and its checks.

\subsection{Cosmology without a new source law}
The cosmological constant in \eqref{eq:SeffLambda} may be moved to the matter side as
\begin{equation}
T^{(\Lambda)}_{\mu\nu}=-\frac{\Lambda c^4}{8\pi G}g_{\mu\nu}.
\end{equation}
That specific constant vacuum tensor is equivalent to $\Lambda$. A general dark-energy field or fluid is not; its dynamics and perturbations must be supplied separately. The same component must not be counted on both sides of the field equation.

For a Friedmann--Lema\^itre--Robertson--Walker (FLRW) model with non-$\Lambda$ density $\rho$ and pressure $p$,
\begin{equation}
\frac{\ddot a}{a}=-\frac{4\pi G}{3}\left(\rho+\frac{3p}{c^2}\right)
+\frac{\Lambda c^2}{3}.
\end{equation}
Positive $\Lambda$ contributes toward acceleration but need not outweigh matter at every epoch. The supernova and Planck references provide motivation for studying accelerated cosmological benchmarks \parencite{Riess1998,Perlmutter1999,Planck2018Parameters}. They do not identify dark energy with $\PhiSrc$, derive the value of $\Lambda$, or establish a source origin of GR.

\section{What an observer-based flow description can explain}
Once a spacetime and a family of observers have been selected, the word flow has a precise mathematical use. A timelike congruence has four-velocity $u^\mu=\dd x^\mu/\dd\tau$, normalized by $u^\mu u_\mu=-c^2$. Its local rest-space projector is
\begin{equation}
h_{\mu\nu}=g_{\mu\nu}+\frac{u_\mu u_\nu}{c^2}.
\end{equation}
The covariant derivative separates as
\begin{equation}
\nabla_\mu u_\nu=\frac13\theta h_{\mu\nu}+\sigma_{\mu\nu}
+\omega_{\mu\nu}-\frac{u_\mu a_\nu}{c^2}.
\end{equation}
Here $\theta$ is the local volume-change rate, $\sigma$ is the anisotropic deformation rate, $\omega$ is vorticity, and $a$ is proper acceleration. This is the standard congruence description \parencite{EllisVanElst1999}.

The description is useful because it separates effects that a single image of moving substance would merge. In homogeneous cosmology, the comoving congruence has $\theta=3H$ and zero shear, vorticity, and acceleration. For static supported observers, all three spatial kinematic rates can vanish while the proper acceleration is $a_\mu=c^2D_\mu\ln N$. An expanding universe therefore need not swirl, and static gravity need not involve expanding observer volume.

Rotation requires additional care. The zero-angular-momentum observers used in a circular stationary-axisymmetric geometry are hypersurface normal and have zero vorticity, even though their angular velocity relative to an appropriate asymptotic reference can reflect frame dragging \parencite{Gourgoulhon2010}. Conversely, rotating observers in flat spacetime can have nonzero vorticity. The congruence and reference observers must be specified before assigning a physical interpretation to rotational quantities.

Likewise, shear is not curvature itself. In a weak vacuum transverse--traceless (TT) wave, a chosen freely falling congruence has $\sigma_{ij}=\dot h^{\TT}_{ij}/2$ at first order, while its tidal acceleration per unit separation involves $\ddot h^{\TT}_{ij}/2$. A static congruence can have zero shear in the presence of nonzero tides. Geometry gives these relations, the curvature convention, and the corrected Raychaudhuri identity in full.

The result is an effective interpretive vocabulary, not a source mechanism. All these quantities are calculated after $g$ and $u$ have been supplied. Declaring them to be manifestations of $\PhiSrc$ without a map would rename target objects rather than derive them.

\section{Explanatory value and empirical limits}
A physical interpretation may be useful by clarifying concepts, suggesting a mathematical construction, or connecting otherwise separate phenomena. Those possibilities should be evaluated directly. They do not follow automatically from exact agreement with a theory that was adopted at the outset.

For this proposal, the key explanatory question is whether source assumptions can eventually reduce what must be independently assumed at effective level. If a derivation simply postulates a Lorentzian metric with Einstein--Hilbert dynamics under new names, the measured theory may be sound but the intended explanation has not advanced. If instead independently specified source relations imply an operational causal structure, a metric, or part of its dynamics, the scope and assumptions of that result can be assessed.

Empirical equivalence also limits confirmation. Where the proposed description and ordinary GR assign the same outcomes to every measurement by construction, a successful GR test does not discriminate in favor of this source ontology over GR without it. This does not logically refute the ontology. It means that its additional explanatory claim needs evidence of its own, such as a demonstrable reduction of assumptions or a discriminating prediction from a completed model.

The sparse source specification cannot currently support a full consistency assessment. Without defined source states and dynamics, one cannot count its propagating modes, derive its conservation laws, or establish its stability. The absence of those definitions is a concrete limitation, not a theorem that no suitable definitions can be found.

\section{Relationship to other gravitational approaches}
\subsection{Standard GR and preferred-structure theories}
The effective benchmark is ordinary GR. Congruence kinematics and Arnowitt--Deser--Misner (ADM) variables are standard ways to analyze that theory, not additional evidence for a source ontology. The word \Aether{} must therefore be distinguished from its use in other actions.

Einstein-aether theory introduces a dynamical unit timelike vector and studies its coupled gravitational modes \parencite{JacobsonMattingly2001,JacobsonMattingly2004}. Ho\v{r}ava's proposal instead uses anisotropic short-distance scaling and a different spacetime-symmetry setup \parencite{Horava2009}. The present effective action introduces neither structure. A future source completion that produced observable preferred structure would need a new analysis rather than relying on this terminological distinction.

\subsection{Derivational and analogue approaches}
Several established lines of work seek relations between Einsteinian dynamics and other principles. Jacobson's thermodynamic argument uses horizon entropy and the heat relation for local causal horizons \parencite{Jacobson1995}. Induced-gravity arguments concern how gravitational terms can arise in an effective description of other fields \parencite{Visser2002}. Self-coupling arguments begin with a defined gauge field and use consistency requirements to motivate nonlinear interactions \parencite{Deser2004}. Analogue-gravity systems can reproduce aspects of effective Lorentzian propagation without automatically reproducing the full Einstein dynamics \parencite{BarceloLiberatiVisser2011}.

These are comparisons of explanatory tasks, not derivations of the present proposal. Each begins with definite mathematical and physical assumptions. The untyped source placeholder here does not inherit their conclusions. Related questions about canonical reconstruction and uniqueness likewise require their own hypotheses; general covariance or four-dimensionality alone is not the missing argument.

\section{A dual physics and AI research objective}
The project's physical aim is to determine whether a specified development of the source ontology can account for relativistic structure without importing the claimed result as a premise. A constructive derivation is one possible outcome. A rigorous obstruction within a declared model class is another. The program does not require a favorable verdict on the ontology to count the investigation as scientifically useful.

The second aim is to study AI-assisted research itself. A system must do more than generate a plausible account of gravity. It must formulate the question, keep assumptions distinct from consequences, check its calculations, seek counterexamples, and report uncertainty when an argument remains incomplete. A negative physical result can demonstrate successful reasoning; an attractive but invalid derivation cannot.

Sys4AI is the intended agentic AI framework for organizing research roles, tasks, tool use, and verification records. These are proposed uses in this project, not reported implementation results. The following protocol defines a research case study that could be implemented and independently assessed. It does not establish the performance of any particular model, agent arrangement, or software release.

\section{Task specification and the meaning of success}
\subsection{Freeze the question before assessing the answer}
A task must identify a source version, its mathematical objects and laws, the admissible source models, and the permitted reconstruction maps. It must name the effective target and its physical domain. Examples of different targets are a local causal structure, an operational clock, gravitational dynamics in a controlled limit, and exact recovery of a specified classical GR model. A result about one target must not be promoted to another by changing the wording after the search.

The source assumptions need not be complete before AI participates. In a fixed-assumption audit, an agent analyzes the supplied proposal without adding premises. In open construction, it may propose missing definitions or laws, but must label every addition and explain its role. Changes create a new version of the question. The original and strengthened assumption sets must remain distinguishable in the record.

A task must also address existence. An inconsistent source specification does not derive physically meaningful GR through a vacuous implication. An admissible example, a suitable existence argument, or an explicit unresolved existence condition is needed. A theorem may be conditional on such a condition, but its research status must state that dependence.

The protocol does not demand that every task derive all matter or every measured constant. A gravitational result may assume a particular matter sector, and an approximate result may apply only at certain scales. Those restrictions belong in the claim. The question is what explanatory content the source construction adds after its independently supplied inputs are accounted for.

\subsection{Separate the outcome categories}
A \emph{verified derivation} establishes the declared implication under the stated assumptions and regime. A proposed derivation remains provisional until the relevant checks support it. Matching one familiar spacetime or one finite set of observables establishes only that narrower result unless a broader argument is supplied.

A \emph{counterexample to entailment} satisfies the task's source assumptions but violates its required conclusion under the specified reconstruction. It disproves that implication. It does not establish that every other source completion or reconstruction must fail. In particular, premises may permit both GR-compatible and non-GR completions, in which case they do not select GR uniquely.

A \emph{scoped obstruction} establishes that no admissible member of the defined model or reconstruction class can meet the target under the stated conditions. Its quantifiers matter: excluding one construction is not excluding every possible construction. A conclusion about a version of the ontology does not automatically apply to later versions with different premises. A hypothetical model forcing all reconstructed metrics to be flat would exclude recovery of a target class containing curved solutions, but this restriction has not been proved for the present ontology.

An \emph{unresolved search} produces neither a verified positive result nor a valid negative argument. It may yield useful lemmas, calculations, or a precise missing step, which should be reported separately. Exhausting time, tokens, or a list of candidate models is not a proof of impossibility unless the task itself is a justified finite exhaustive decision problem and the search is complete.

Finally, a construction showing that some permitted completion is compatible with GR is an existence result at that scope. It is not a proof that the original assumptions force GR, nor that the completion is unique. The evaluation must preserve this distinction between compatibility, entailment, and exclusion.

\section{Proposed evaluation protocol}
\subsection{Evidence beyond restating the manuscripts}
The existing papers already document the lack of a source-to-GR reconstruction. An agent that repeats that finding has performed an audit or summary, not demonstrated new research. Credit for research progress requires an additional inspectable result: a precise source specification, a nontrivial implication, a justified reconstruction step, a counterexample, or a previously unrecognized defect in an argument.

Assessment should record the physical outcome separately from the correctness and scope of the research artifact. A correct negative theorem should not be scored as a reasoning failure merely because it rules out the candidate. A partial derivation should receive credit for its proved portion without being called full recovery. Unsupported confidence, concealed assumptions, and invalid counterexamples should be identified directly rather than hidden in an aggregate score.

For each result, retain the claim statement, definitions, assumptions, dependencies, derivation or witness, and checking record. Distinguish calculations performed from checks merely proposed. A missing premise cannot be repaired by a convincing narrative. The final claim must be readable and checkable without access to unrecorded agent exchanges.

\subsection{Independent checking and physical interpretation}
A checking stage should seek errors rather than only confirm the proposed result. It should inspect assumptions, intermediate steps, limiting cases, and possible counterexamples. Where practical, it should use a separate calculation, implementation, or reviewer. Assigning an agent a critic role is a process choice; it is not itself evidence that the result was independently validated.

Symbolic computation can check specified algebraic steps. Numerical calculations can test stated cases with reported precision and error control. Formal proof tools can verify a precisely encoded theorem. None of these checks alone establishes that the encoding or model captures the intended physics. Qualified human assessment is needed for the physical interpretation and scope of any consequential recovery or no-go claim. Unchecked parts remain visible, and agreement among agents is not substituted for a proof.

The verification record should identify who or what checked each component, which inputs were used, and the limits of the check. No claim of external review or human approval follows merely from including that requirement in the protocol.

\subsection{Controlled comparisons and reproducible records}
Before comparing models, fix the task set, source snapshot, permitted information, tool access, resource limits, and assessment criteria. Record the model identifier and available version information, settings, prompts, tool and code outputs, elapsed time, token use, actual costs where available, and human interventions. Record seeds where tools expose them; a seed alone does not guarantee identical reruns. Missing resource measurements must be labeled unavailable rather than estimated as if observed.

The comparison unit must be explicit. A base model with a fixed prompt is different from a system using several agents, retrieval, memory, symbolic tools, and human guidance. A result from the full Sys4AI arrangement demonstrates the behavior of that arrangement under the recorded conditions. It does not isolate the base model's contribution.

Useful comparisons would hold the model fixed while varying the workflow, or hold the workflow and task budget fixed while varying the model. A direct single-agent workflow provides a simple baseline. Removing a component in a controlled comparison can test whether it contributes; such a test is often called an ablation. Human help should either be controlled across conditions or reported as a separate condition. Repeated trials and their variability are preferable to selecting a favorable run.

\subsection{Controls, exposure, and the limit of generalization}
The inherited GR calculations can serve as known-answer control tasks. Carefully specified flawed derivations can test error detection; known counterexamples can test whether the system distinguishes non-entailment from a class-wide obstruction. Such controls should have independently checked answers and criteria established before model outputs are graded. Their usefulness must be tested, not assumed from this proposal.

Evaluation should distinguish access to source physics from access to evaluation answers. The manuscripts may be legitimate task inputs, but repeating an answer supplied in them cannot demonstrate its independent discovery. Record which documents and retrieval sources were available. Keep assessment-only material and held-out checks separate where the task design permits, and report uncertainty about prior model exposure rather than claiming it can always be excluded.

This single ontology is a research case study. Even a strong checked result would support a claim about the demonstrated task, not general scientific competence. Broader evaluation would require additional problems, stable assessment criteria, repeated trials, and examination of failure cases. No pass rate, model ranking, cost comparison, or validated benchmark score is available in the present manuscripts.

\section{Initial feasibility study and stopping conditions}
The first study should select one explicit development of the source proposal and one attainable mathematical question. A causal-geometry or clock-reconstruction target is a possible starting point, not a source assumption adopted here. The study should expose its premises, give a construction or admissible example where possible, and attempt a checkable implication or obstruction.

A positive result justifies only the next task supported by its dependencies. For example, a metric construction may motivate a separate dynamics task; it does not complete that task in advance. A negative result should identify the exact restriction responsible and the alternatives it leaves open. An unresolved result should locate the missing step and preserve the failed attempts needed to avoid repeating them without new information.

Stop a run when its predefined resource limit is reached, when the declared result has been checked to the agreed standard, or when further work would change the question without authorization. New assumptions or a larger target belong in a new task version. The research deliverable is a claim with evidence and scope, not a requirement to reach either a positive or negative verdict at any cost.

\section{Conclusion}
The \Aflow{} Interpretation of Relativity presently combines a tentative source ontology with an adopted GR model. Observer congruences make its effective flow language precise, but do not convert that language into a derivation of spacetime, matter, or clocks. Its standard relativistic predictions are inherited under explicit matching conditions.

The proposal's unresolved question is therefore clear: can a defined source structure produce some or all of the effective relational and dynamical content without importing that content as a premise? The adopted GR model supplies an intelligible benchmark for addressing that question. This benchmark neither establishes the physical reality of the proposed source nor shows that all future source constructions must fail.

The dual research objective makes both constructive and negative results relevant, but it does not make either available by declaration. A useful first contribution is a checkable result about a defined source model, accompanied by a reproducible research record. A result against that model can still demonstrate capable reasoning. The project remains an open physical proposal and a proposed AI research case study until such evidence is produced.

\References
\end{document}
